\documentclass[10pt,a4paper]{amsart}
\usepackage[margin=19mm]{geometry}
\usepackage{amsmath,amssymb,mathtools}
\usepackage[scr=rsfs,cal=euler]{mathalfa}
\usepackage{xfrac}
\usepackage[unicode,hidelinks,bookmarks=false]{hyperref}
\newcommand{\ie}{i.e.\ }
\newcommand{\cf}{cf.\ }
\newcommand{\resp}{respectively\ }
\newcommand{\Subsection}{\subsection}
\providecommand{\nobreakdash}{\nobreak\hbox{-}}
\setlength{\emergencystretch}{2em}
\multlinegap0pt
\usepackage{amssymb}
\newtheorem{definition}{Definition}
\newtheorem{lemma}{Lemma}
\title{Spectral asymptotics for two-term even-order differential operators with homogeneous delay}
\author{Dmitry M. Polyakov}
\date{Source: arXiv:2606.21263v1 (CC BY 4.0)}
\begin{document}
\maketitle

\noindent\emph{Source and licence.} Spectral asymptotics for two-term even-order differential operators with homogeneous delay. Version arXiv:2606.21263v1 (CC BY 4.0), distributed by arXiv under the Creative Commons Attribution 4.0 International licence, http://creativecommons.org/licenses/by/4.0/, as recorded in the arXiv licence metadata for this exact version and shown on the abstract page https://arxiv.org/abs/2606.21263v1. Full licence notice: \texttt{licenses/arxiv-2606.21263-CC-BY-4.0.txt}.\par\smallskip

\noindent\textit{Editorial scope.} Counted unit: the article's lemma lh1 (source0.tex lines 629--631). The article's complete original proof of it is delivered with it (source0.tex lines 632--678, one \texttt{proof} environment of 47 lines). No mathematics is written, repaired or re-derived: the statement and the proof are the article's own lines, in the article's own notation, and no retained line is substituted or re-flowed.  Retained uncounted as the source-local context the proof needs: lines 515--547; lines 528--530; lines 583--586; lines 588--592; lines 609--612; lines 613--617; lines 620--623.  Nothing was omitted from any retained span, so the package carries no omission record.  The retained lines cite 1 works, whose entries are transcribed verbatim from the article's own bibliography source in the e-print and delivered with short editorial labels recorded in the item's reference mapping.  No retained line contains a figure environment, an image-inclusion command or drawing code, so no image is delivered and the package declares no asset dependency.  Editorial formatting only, in addition: an AMS \texttt{amsart} preamble that restates the article's own declarations and macros used by the retained lines, namely the environments \texttt{definition}, \texttt{lemma} on separate counters, together with the standard packages those lines need.  The retained environments print in this document as Definition 1, Lemma 1 in order of appearance, whereas the article prints the same items with the numbering of its own full document.  The complete native source of the article as distributed in the arXiv e-print for this exact version is bundled unchanged as \texttt{sources/203262/source0.tex}.
\begin{definition}[\cite{BaskKrUsk}]\label{defadm}
Suppose that $A$ is a closed, densely defined linear operator, $\mathfrak{U}$
is a linear subspace of $\mathfrak{L}_A(H)$, $J:\mathfrak{U}\to\mathfrak{U}$,
and $\Gamma:\mathfrak{U}\to B(H)$. The collection $(\mathfrak{U}, J, \Gamma)$ is an
admissible triplet for the operator $A$ and $\mathfrak{U}$ is the space of
admissible perturbations, if the following properties hold.

1) The space $\mathfrak{U}$ is a Banach space (with respect to some norm
$\|\cdot\|_{\mathfrak{U}}$) that is continuously embedded in $\mathfrak{L}_A(H)$.

2) $J$ and $\Gamma$ are continuous linear operators and, moreover, $J$ is a projection.

3) $(\Gamma X)D(A)\subset D(A)$ and
\begin{equation}\label{adA}
A(\Gamma X)x-(\Gamma X)Ax=(X-JX)x
\end{equation}
for all $x\in D(A)$, $X\in\mathfrak{U}$. Moreover, $Y=\Gamma X\in B(H)$
is the unique solution of the equation
\[
AY-YA=X-JX,
\]
that satisfies $JY=0$.

4) $X\Gamma Y$, $(\Gamma X)Y\in\mathfrak{U}$ for all $X$, $Y\in\mathfrak{U}$,
and there is a constant $\gamma>0$ such that
\[
\|\Gamma\|\leqslant \gamma, \qquad \max\{\|X\Gamma Y\|_{\mathfrak{U}},
\|(\Gamma X)Y\|_{\mathfrak{U}}\}\leqslant \gamma\|X\|_{\mathfrak{U}}\|Y\|_{\mathfrak{U}}.
\]

5) For every $X\in\mathfrak{U}$ and $\delta>0$ there exists a number
$\lambda_\delta\in\rho(A)$ such that~$\|X(A-\lambda_\delta I)^{-1}\|<\delta$.
\end{definition}

\begin{equation}
A(\Gamma X)x-(\Gamma X)Ax=(X-JX)x
\end{equation}

\begin{equation}\label{2.0}
\mu_n:=(\pi n)^{2k},\quad e_n(x):=\sqrt{2}\sin\pi n x,
\quad \mathcal{P}_nx:= (x, e_n)_{L^2(0, 1)}e_n, \quad x\in [0, 1], \quad n\geqslant 1,
\end{equation}

\begin{equation}\label{3.0}
\mu_n:=(\pi n)^{2k},\quad e_n(x):=\sqrt{2}\cos\pi nx,
\quad \mathcal{P}_nx:= (x, e_n)_{L^2(0, 1)}e_n, \quad x\in [0, 1],\quad
n\geqslant 0,
\end{equation}

\begin{equation}\label{j}
(J\mathcal{X})_{lj}=\delta_{l-j}\mathcal{X}_{lj}, \quad l, j\in\mathbb{J}, \quad
\mathcal{X}\in\mathfrak{U},
\end{equation}

\begin{equation}\label{gamma}
(\Gamma\mathcal{X})_{lj}=\frac{\mathcal{X}_{lj}}{\mu_l-\mu_j},
\qquad l\ne j, \qquad (\Gamma\mathcal{X})_{ll}=0, \qquad l, j\in\mathbb{J},
\qquad \mathcal{X}\in\mathfrak{U}.
\end{equation}

\begin{equation}\label{estJG}
\|J\|_{B(L^2(0, 1))}=1, \qquad \|\Gamma\|_{B(L^2(0, 1))}\leqslant\frac{\pi}{2\gamma_n}\leqslant\frac{2}{\gamma_n},
\qquad \gamma_n=\min_{l\ne n\in\mathbb{J}}|\mu_n-\mu_l|.
\end{equation}

\begin{lemma}\label{lh1}
The triplet $(\mathfrak{U}, J, \Gamma)$ is admissible.
\end{lemma}

\begin{proof}
Property~1) of Definition~\ref{defadm} is immediate from the definitions of
$\mathfrak{U}=B(L^2(0, 1))$. Property 2) follows from the definitions of $J$ and $\Gamma$
and \cite[Lemma~3.4]{BaskKrUsk}.

Now we prove Property~3). The proof is identical for both the Dirichlet and Neumann type boundary conditions.
Let $x\in D(\mathcal{A})$ and $\lambda\in\rho(\mathcal{A})$. There exists
$y\in L^2(0, 1)$ such that
\[
x=(\mathcal{A}-\lambda I)^{-1}y=\sum_{j\in\mathbb{J}}\frac{1}{\mu_j-\lambda}\mathcal{P}_jy,
\]
where $\mathcal{P}_j$, $j\in\mathbb{J}$, are defined by \eqref{2.0} and \eqref{3.0}. Then
for $\mathcal{X}\in\mathfrak{U}$ we get
\begin{align*}
(\Gamma\mathcal{X})x &= (\Gamma\mathcal{X})(\mathcal{A}-\lambda I)^{-1}y=
\sum\limits_{\substack{l, j\in\mathbb{J} \\ l\ne j}}
\frac{\mathcal{P}_l\mathcal{X}\mathcal{P}_jy}{(\mu_l-\mu_j)(\mu_j-\lambda)} \\
&= \sum\limits_{\substack{l, j\in\mathbb{J} \\ l\ne j}}
\frac{\mathcal{P}_l\mathcal{X}\mathcal{P}_jy}{(\mu_l-\mu_j)(\mu_l-\lambda)}+
\sum\limits_{\substack{l, j\in\mathbb{J} \\ l\ne j}}
\frac{\mathcal{P}_l\mathcal{X}\mathcal{P}_jy}{(\mu_l-\lambda)(\mu_j-\lambda)} \\
&= (\mathcal{A}-\lambda I)^{-1}(\Gamma\mathcal{X})y+(\mathcal{A}-\lambda I)^{-1}
(\Gamma\mathcal{X})(\mathcal{A}-\lambda I)^{-1}y \\
&= (\mathcal{A}-\lambda I)^{-1}(\Gamma\mathcal{X})(x+y)\in D(\mathcal{A}).
\end{align*}
Thus, $(\Gamma\mathcal{X})D(\mathcal{A})\subset D(\mathcal{A})$. Moreover,
these identities show that the matrices of the bounded operators
$(\Gamma\mathcal{X})(\mathcal{A}-\lambda I)^{-1}
-(\mathcal{A}-\lambda I)^{-1}(\Gamma\mathcal{X})$ and $(\mathcal{A}-\lambda
I)^{-1}(\mathcal{X}-J\mathcal{X})(\mathcal{A}-\lambda I)^{-1}$ coincide.
This gives the formula \eqref{adA}. It follows from \eqref{j} and \eqref{gamma}
that the identity $J(\Gamma\mathcal{X})=0$ holds. Therefore, Property 3) is proved.

Now we verify Property~4). Assume that $\mathcal{X}$, $\mathcal{Y}\in B(L^2(0, 1))$.
Using the estimates \eqref{estJG}, we obtain
\[
\|\mathcal{X}\Gamma\mathcal{Y}\|_{B(L^2(0, 1))}\leqslant \|\mathcal{X}\|_{B(L^2(0, 1))}
\|\Gamma\|_{B(L^2(0, 1))}\|\mathcal{Y}\|_{B(L^2(0, 1))}\leqslant
2\gamma_n^{-1}\|\mathcal{X}\|_{B(L^2(0, 1))}\|\mathcal{Y}\|_{B(L^2(0, 1))}.
\]
Obviously, this inequality holds for $\|(\Gamma\mathcal{X})\mathcal{Y}\|_{B(L^2(0, 1))}$.
Property 4) is proved with constant $2\gamma_n^{-1}$.

Finally, Property~5) follows directly from the boundedness of the operator
$\mathcal{X}(\mathcal{A}-\lambda I)^{-1}$ (see also~\cite[Definition~2.1]{BaskKrUsk}).
Lemma is proved.
\end{proof}

\begin{thebibliography}{99}
\bibitem[BaskKr]{BaskKrUsk} A.G. Baskakov, I.A. Krishtal, N.B. Uskova. Similarity techniques in the spectral analysis of perturbed operator matrices. J. Math. Anal. Appl. 2019. V.~477. P.~930--960.
\end{thebibliography}
\end{document}
